% Volume VIII: Computational Neuroscience Mathematics
% Continuity note for Chapter 48.
% Previous dependencies: Chapter 5's projections and least squares; Chapter 6's eigenvectors, SVD, covariance, and PCA; Chapter 7's tensor shapes; Chapter 13's random vectors and covariance matrices; Chapter 15's conditional expectation; Chapter 25's dynamical systems; Chapter 27's filtering concepts when available; Chapter 47's population response vectors and decoders.
% New durable concepts: neural state spaces; activity matrices; neural trajectories; covariance geometry; PCA for neural data; factor models; latent state-space models; Kalman prediction-update intuition; neural manifolds; intrinsic dimension; tangent spaces; behavioral readouts from latent coordinates.
% Forward hooks: Chapter 49 studies recurrent dynamics that can generate manifolds, attractors, and oscillations; Chapter 50 treats latent causes and beliefs in Bayesian perception; Chapter 51 uses these models for neural data analysis and model comparison; Chapter 53 connects PCA to neural manifolds in an integrated case study.
% Notation commitments: Population activity at time t is x_t in \mathbb R^N; data matrices are X in \mathbb R^{M\times N} with rows as observations and columns as neurons; centered data are X_c; latent states are z_t in \mathbb R^d with d\ll N; observation matrices are C; process noise is \epsilon_t and observation noise is \eta_t.
\section{Chapter 48: Population Geometry and Latent Neural Dynamics}
\label{ch:population-geometry-latent-neural-dynamics}

Chapter 47 treated a neural population response as a vector that can encode and decode variables. Chapter 48 asks what happens when we stop focusing on a single stimulus and instead study the shape and motion of population activity itself. A recording from \(N\) neurons is a point in \(\mathbb R^N\) at each time or trial. Across time, those points trace trajectories. Across conditions, they form clouds, curves, surfaces, or more complicated low-dimensional sets.

The central problem is that \(N\) can be large, while the meaningful degrees of freedom may be much smaller. A motor task may recruit hundreds of neurons but evolve along a few dominant directions. A working-memory task may keep activity near a low-dimensional attractor. A decision task may compress population activity onto a latent decision variable. The mathematics of geometry, projections, latent variables, and dynamical systems makes these claims precise.

\paragraph{Chapter problem.}
How can high-dimensional neural recordings be represented as state-space trajectories, low-dimensional projections, latent dynamical systems, and manifolds linked to behavior?

\paragraph{Section map.}
Section 48.1 defines neural state spaces and activity matrices. Section 48.2 derives PCA and connects it to factor models for dimensionality reduction. Section 48.3 introduces latent dynamical systems with state and observation equations. Section 48.4 defines neural manifolds and shows how low-dimensional coordinates can support behavioral readouts.

\subsection{48.1 Neural state spaces}

\paragraph{Motivating question.}
How can simultaneous activity from many neurons be viewed as a point or trajectory in a geometric space?

\paragraph{Beginner setup.}
At one time bin, a population of \(N\) neurons produces \(N\) numbers. These may be spike counts, firing-rate estimates, calcium fluorescence values, or latent firing rates. Put them into a vector
\[
x_t=(x_{t1},\ldots,x_{tN})\in\mathbb R^N.
\]
This vector is the neural state at time \(t\). If we record across \(T\) time bins, the sequence
\[
x_1,x_2,\ldots,x_T
\]
is a trajectory in \(N\)-dimensional neural state space.

A trial-by-time-by-neuron data tensor \(R_{b,t,n}\) can be reshaped into a data matrix
\[
X\in\mathbb R^{M\times N},
\]
where each row is one observation, such as one trial-time pair, and each column is one neuron. Geometry enters because rows can be compared by distances, angles, projections, and covariance.

\paragraph{Formal setup.}
Let \(N\) be the number of recorded neurons. A neural state space is the vector space \(\mathbb R^N\) equipped with coordinates corresponding to neurons or features. An observation at index \(m\) is
\[
x_m\in\mathbb R^N.
\]
Stacking \(M\) observations gives
\[
X=
\begin{bmatrix}
x_1^\top\\
\vdots\\
x_M^\top
\end{bmatrix}
\in\mathbb R^{M\times N}.
\]
The empirical mean activity is
\[
\bar x=\frac1M\sum_{m=1}^M x_m\in\mathbb R^N,
\]
and the centered data matrix is
\[
X_c=
\begin{bmatrix}
(x_1-\bar x)^\top\\
\vdots\\
(x_M-\bar x)^\top
\end{bmatrix}.
\]
The empirical covariance matrix is
\[
\boxed{
C=\frac{1}{M-1}X_c^\top X_c\in\mathbb R^{N\times N}.
}
\]
Its diagonal entries are variances of individual neurons; its off-diagonal entries are covariances between neurons.

\paragraph{Core intuition.}
A state-space plot is not a new measurement. It is a geometric view of the same population activity. If \(N=2\), each time point is a point in a plane. If \(N=3\), it is a point in 3D. For \(N=100\), we cannot draw the space directly, but the same geometry exists algebraically.

Trajectories matter because neural activity is not just a cloud of unrelated points. During movement, a population may rotate through state space. During a decision, it may drift toward one of two regions. During memory, it may stay near a stable low-dimensional set. The geometry tells us what changes, what stays stable, and which directions carry variance or task information.

\paragraph{Main mechanism: distances, covariance, and trajectories.}
The Euclidean distance between two neural states is
\[
\|x_a-x_b\|_2
=\left(\sum_{i=1}^N (x_{ai}-x_{bi})^2\right)^{1/2}.
\]
The squared distance from the mean, averaged over observations, is the total variance:
\[
\frac{1}{M-1}\sum_{m=1}^M \|x_m-\bar x\|_2^2
=\operatorname{tr}(C).
\]
To see this, expand
\[
\sum_m \|x_m-\bar x\|_2^2
=\sum_m\sum_i (x_{mi}-\bar x_i)^2.
\]
The right side is \((M-1)\) times the sum of the covariance diagonal entries, which is \((M-1)\operatorname{tr}(C)\).

Thus covariance geometry is not an optional add-on. It summarizes how the population cloud spreads around its mean, and the trace gives the total variance available for dimensionality reduction.

\paragraph{Worked examples.}
\emph{Small numerical example.}
Suppose three time bins from two neurons give
\[
x_1=(1,0),\qquad x_2=(2,1),\qquad x_3=(3,1).
\]
The activity moves from low neuron-1 activity to high neuron-1 activity while neuron 2 rises and then stays high. The mean is
\[
\bar x=(2,2/3).
\]
The distance between \(x_1\) and \(x_3\) is
\[
\|(1,0)-(3,1)\|_2=\sqrt{5}.
\]
This sequence is a short trajectory in a two-neuron state space.

\emph{Computational neuroscience example.}
In a reaching experiment, each time bin during a movement can be represented by a vector of motor-cortex firing rates. Plotting the top two or three projected dimensions often reveals rotational or preparatory trajectories that are difficult to see neuron by neuron.

\emph{AI example.}
Hidden states of a recurrent neural network or transformer residual stream can be treated as state-space points. Across time or layers, these points trace trajectories through representation space. The same tools, such as covariance, projections, and trajectory distances, can compare artificial and biological representations.

\paragraph{Common misconceptions and failure modes.}
A neural state space is not necessarily physical space in the brain. Its axes are measured variables or features, not anatomical dimensions. Euclidean distance is not automatically the right similarity measure; scaling, noise variance, and preprocessing change distances. A low-dimensional plot is a projection, not the full state space. Trial averaging can reveal reliable trajectories but can also hide meaningful trial-to-trial variability.

\paragraph{Exercises.}
\begin{enumerate}[label=\textbf{48.1.\arabic*.}, leftmargin=*]
\item \textbf{Mechanical.} A recording has \(40\) trials, \(100\) time bins, and \(60\) neurons. If each trial-time pair is one observation, what are \(M\), \(N\), and the shape of \(X\)?
\item \textbf{Proof or derivation.} Prove that \(\operatorname{tr}(C)\) equals the sum of empirical variances across neurons.
\item \textbf{Numerical/computational.} For \(x_1=(0,1,2)\) and \(x_2=(2,1,0)\), compute their Euclidean distance. Then compute their dot product.
\item \textbf{Interpretation.} What does it mean for two task conditions to form separated clouds in population state space?
\item \textbf{Debug the misconception.} A student says, "A 3D PCA plot is the neural state space." Correct the statement by distinguishing the full state space from a projection.
\end{enumerate}

\subsection{48.2 Dimensionality reduction for neural data}

\paragraph{Motivating question.}
How can high-dimensional neural activity be summarized by a few coordinates while preserving the most important structure?

\paragraph{Beginner setup.}
Dimensionality reduction looks for a simpler coordinate system. In neural data, the number of neurons \(N\) may be large, but the activity may vary mostly along a few directions. Principal component analysis, or PCA, finds orthogonal directions of largest variance. Factor analysis and related latent-variable models add explicit noise assumptions.

The simplest picture is a cloud of points in \(\mathbb R^N\) stretched along one direction. Projecting onto that direction captures most variation. Projecting onto an orthogonal low-variance direction captures little. PCA makes this precise using the covariance matrix.

\paragraph{Formal setup.}
Let \(X_c\in\mathbb R^{M\times N}\) be centered data and
\[
C=\frac{1}{M-1}X_c^\top X_c.
\]
Because \(C\) is symmetric positive semidefinite, it has orthonormal eigenvectors
\[
v_1,\ldots,v_N
\]
with eigenvalues
\[
\lambda_1\geq \lambda_2\geq\cdots\geq \lambda_N\geq0.
\]
The \(k\)-dimensional PCA projection matrix is
\[
V_k=[v_1\ \cdots\ v_k]\in\mathbb R^{N\times k}.
\]
The low-dimensional coordinates are
\[
\boxed{
Z=X_cV_k\in\mathbb R^{M\times k}.
}
\]
The rank-\(k\) reconstruction in the original neural coordinates is
\[
\widehat X_c=ZV_k^\top=X_cV_kV_k^\top.
\]
The fraction of variance explained by the first \(k\) components is
\[
\boxed{
\frac{\lambda_1+\cdots+\lambda_k}{\lambda_1+\cdots+\lambda_N}.
}
\]

\paragraph{Core intuition.}
PCA rotates the coordinate system so that the first axis points along the largest spread of the data, the second axis points along the largest remaining orthogonal spread, and so on. Projection onto the first few axes gives a low-dimensional summary.

Factor analysis changes the modeling story. It assumes a latent variable \(z_m\in\mathbb R^k\) generates the observation:
\[
x_m=Cz_m+\eta_m,
\]
where \(C\in\mathbb R^{N\times k}\) is a loading matrix and \(\eta_m\) is observation noise. PCA is geometric and variance based; factor analysis is probabilistic and separates shared latent variability from private noise under its assumptions.

\paragraph{Main theorem: PCA minimizes linear reconstruction error.}
Among all \(k\)-dimensional orthonormal projection matrices \(VV^\top\), with \(V\in\mathbb R^{N\times k}\) and \(V^\top V=I_k\), PCA minimizes
\[
\|X_c-X_cVV^\top\|_F^2.
\]
The minimum is achieved by \(V=V_k\), the top \(k\) eigenvectors of \(C\).

\emph{Proof sketch.}
The reconstruction error decomposes as
\[
\|X_c-X_cVV^\top\|_F^2
=\|X_c\|_F^2-\|X_cV\|_F^2.
\]
The first term is fixed, so minimizing error is the same as maximizing projected variance:
\[
\|X_cV\|_F^2
=\operatorname{tr}(V^\top X_c^\top X_c V)
=(M-1)\operatorname{tr}(V^\top C V).
\]
The Rayleigh-Ritz theorem says this trace is maximized by choosing the columns of \(V\) to be the eigenvectors with largest eigenvalues. Therefore top-PC projection gives the best rank-\(k\) orthogonal linear reconstruction.

\paragraph{Worked examples.}
\emph{Small numerical example.}
Suppose the covariance eigenvalues are
\[
\lambda=(9,4,1,1).
\]
The total variance is \(15\). The first component explains \(9/15=60\%\). The first two explain \((9+4)/15\approx 86.7\%\). A two-dimensional PCA plot therefore preserves most, but not all, variance.

\emph{Computational neuroscience example.}
In population recordings from motor cortex, PCA can reveal that most trial-averaged activity during a task lies in a subspace of much smaller dimension than the number of recorded neurons. This does not mean only a few neurons matter. It means their coordinated activity is well described by a few population directions.

\emph{AI example.}
PCA is used to visualize embeddings from neural networks, compress activations, and diagnose representation collapse. Autoencoders generalize the projection-reconstruction idea by learning nonlinear encoders and decoders, but they need training objectives and validation because nonlinear flexibility can overfit.

\paragraph{Common misconceptions and failure modes.}
PCA preserves variance, not necessarily task information. A low-variance direction can be highly predictive of behavior. A high-variance direction can reflect motion artifact, arousal, or global gain. PCA components are sign-ambiguous: \(v\) and \(-v\) describe the same one-dimensional subspace. The number of dimensions is not determined by a single plot; it depends on the criterion, noise model, task, and validation procedure.

\paragraph{Exercises.}
\begin{enumerate}[label=\textbf{48.2.\arabic*.}, leftmargin=*]
\item \textbf{Mechanical.} If covariance eigenvalues are \(5,3,2\), compute the variance explained by the first one and first two PCs.
\item \textbf{Proof or derivation.} Fill in the algebra showing \(\|X_c-X_cVV^\top\|_F^2=\|X_c\|_F^2-\|X_cV\|_F^2\) when \(V^\top V=I\).
\item \textbf{Numerical/computational.} Given a centered data matrix \(X_c\in\mathbb R^{1000\times 50}\), describe pseudocode to compute the top three PC scores using an eigendecomposition of \(C\).
\item \textbf{Interpretation.} Why might a behaviorally important signal be absent from the first few principal components?
\item \textbf{Debug the misconception.} A student says, "PCA discovers the true latent variables." Correct the statement by distinguishing variance-preserving coordinates from true generative causes.
\end{enumerate}

\subsection{48.3 Latent dynamical systems}

\paragraph{Motivating question.}
How can unobserved low-dimensional states generate high-dimensional neural observations over time?

\paragraph{Beginner setup.}
A latent dynamical system assumes that neural activity is driven by an unobserved state \(z_t\) that changes over time. The state may represent movement phase, decision evidence, internal context, or a compact summary of network dynamics. The observations \(x_t\) are noisy measurements of that state through neurons.

The simplest useful model is linear and Gaussian:
\[
z_{t+1}=Az_t+\epsilon_t,
\qquad
x_t=Cz_t+\eta_t.
\]
The first equation is the state dynamics. The second is the observation model. The latent state dimension \(d\) is smaller than the number of neurons \(N\).

\paragraph{Formal setup.}
Let
\[
z_t\in\mathbb R^d,\qquad x_t\in\mathbb R^N,\qquad d\ll N.
\]
A linear Gaussian state-space model is
\[
\boxed{
z_{t+1}=Az_t+\epsilon_t,\qquad \epsilon_t\sim\mathcal N(0,Q),
}
\]
\[
\boxed{
x_t=Cz_t+\eta_t,\qquad \eta_t\sim\mathcal N(0,R).
}
\]
The matrix \(A\in\mathbb R^{d\times d}\) governs latent dynamics. The loading matrix \(C\in\mathbb R^{N\times d}\) maps latent states to expected neural activity. The covariance \(Q\) is process noise, and \(R\) is observation noise.

The inference problem is to estimate
\[
p(z_t\mid x_1,\ldots,x_t)
\]
or
\[
p(z_{1:T}\mid x_{1:T}).
\]
The first is filtering; the second is smoothing.

\paragraph{Core intuition.}
Dimensionality reduction finds low-dimensional coordinates for observations. A latent dynamical system adds time and a generative story. The latent state moves according to dynamics, then produces noisy observations. Inference reverses this story: from noisy observations, infer the likely latent trajectory.

The Kalman filter is the exact filtering algorithm for the linear Gaussian model. It has a prediction-update rhythm. Prediction uses the dynamics to forecast the next latent state. Update uses the new observation to correct that forecast.

\paragraph{Main mechanism: one-dimensional Kalman update.}
In one dimension, suppose
\[
z_t\mid x_{1:t-1}\sim\mathcal N(m^-,P^-)
\]
is the predicted latent state before seeing \(x_t\), and
\[
x_t=z_t+\eta_t,\qquad \eta_t\sim\mathcal N(0,R).
\]
After observing \(x_t\), the posterior mean is
\[
\boxed{
m=m^-+K(x_t-m^-),
\qquad
K=\frac{P^-}{P^-+R}.
}
\]
The posterior variance is
\[
\boxed{
P=(1-K)P^-.
}
\]
The Kalman gain \(K\) is a precision-weighted trust factor. If observation noise \(R\) is small, \(K\) is close to \(1\), and the estimate moves strongly toward the observation. If prediction uncertainty \(P^-\) is small relative to \(R\), \(K\) is small, and the estimate trusts the prediction more.

\paragraph{Worked examples.}
\emph{Small numerical example.}
Let \(m^-=0\), \(P^-=1\), observation noise \(R=3\), and observed neural summary \(x_t=4\). The Kalman gain is
\[
K=\frac{1}{1+3}=0.25.
\]
The updated mean is
\[
m=0+0.25(4-0)=1.
\]
The posterior variance is
\[
P=(1-0.25)(1)=0.75.
\]
The observation pulls the latent estimate upward, but not all the way to \(4\), because the observation is noisy.

\emph{Computational neuroscience example.}
Latent factor models for reaching data often infer a low-dimensional \(z_t\) whose trajectory evolves smoothly through movement preparation and execution. The observed neurons are noisy readouts of this trajectory. Such models can denoise single trials and reveal dynamics hidden by neuron-wise variability.

\emph{AI example.}
State-space models, recurrent neural networks, and sequence models all maintain hidden states that summarize past observations. A linear Gaussian model is much simpler than a modern sequence model, but it teaches the same distinction between hidden state dynamics and observed outputs.

\paragraph{Common misconceptions and failure modes.}
A latent variable is not directly observed. It is inferred under a model. Rotating the latent coordinates can produce an equivalent model unless constraints fix the coordinate system. Smooth latent trajectories can be artifacts of smoothing assumptions. Linear Gaussian models are useful baselines, but spiking data, nonlinear dynamics, and discrete task events often require extensions.

\paragraph{Exercises.}
\begin{enumerate}[label=\textbf{48.3.\arabic*.}, leftmargin=*]
\item \textbf{Mechanical.} State the shapes of \(A\), \(C\), \(Q\), and \(R\) when \(z_t\in\mathbb R^3\) and \(x_t\in\mathbb R^{80}\).
\item \textbf{Proof or derivation.} Derive the one-dimensional Kalman gain formula by multiplying a Gaussian prior density by a Gaussian likelihood and completing the square.
\item \textbf{Numerical/computational.} Compute the Kalman update for \(m^-=2\), \(P^-=4\), \(R=1\), and \(x_t=5\).
\item \textbf{Interpretation.} What is the difference between filtering \(p(z_t\mid x_{1:t})\) and smoothing \(p(z_t\mid x_{1:T})\)?
\item \textbf{Debug the misconception.} A student says, "The latent trajectory is what the neurons really did." Correct the statement using the distinction between model-inferred state and observed activity.
\end{enumerate}

\subsection{48.4 Neural manifolds and behavior}

\paragraph{Motivating question.}
When high-dimensional neural activity lies near a low-dimensional set, how can that set support behavior and computation?

\paragraph{Beginner setup.}
A neural manifold is a low-dimensional set embedded in high-dimensional neural activity space. If population activity has \(N\) measured coordinates but varies mostly with \(d\) internal coordinates, where \(d\ll N\), then the activity may lie near a \(d\)-dimensional manifold.

For example, a head-direction system can have activity patterns arranged around a ring. The ring is one-dimensional because one angle specifies position on it, but it is embedded in a high-dimensional space of neurons. A reaching movement may trace a two- or three-dimensional trajectory inside a much larger population space.

\paragraph{Formal setup.}
Let \(U\subseteq\mathbb R^d\) be a latent coordinate domain and let
\[
\phi:U\to\mathbb R^N
\]
be a smooth map. The image
\[
\mathcal M=\{\phi(z):z\in U\}
\]
is a \(d\)-dimensional embedded manifold when \(\phi\) is locally one-to-one and its Jacobian
\[
J_\phi(z)\in\mathbb R^{N\times d}
\]
has rank \(d\). The columns of \(J_\phi(z)\) span the tangent space at \(\phi(z)\).

A behavioral readout is a map
\[
b=g(z)
\quad\text{or}\quad
b=h(x),
\]
where \(z\) is a manifold coordinate and \(x\in\mathbb R^N\) is neural activity. If activity is near the manifold, one may first estimate \(z\) from \(x\), then read out behavior.

\paragraph{Core intuition.}
A manifold is more than a pretty low-dimensional plot. It is a claim about degrees of freedom. If a neural population lies near a ring, then one coordinate, angle, explains most of the structured variation. If it lies near a surface, two coordinates may be needed. The manifold can constrain possible activity patterns and make readout easier.

Behavior enters when movement, decision, memory, or perception varies systematically along manifold coordinates. A readout can be simple in latent coordinates even if it is complicated in the original neural coordinates.

\paragraph{Main mechanism: tangent spaces and local decoding.}
Near a point \(z_0\in U\), the manifold map has the local linear approximation
\[
\phi(z_0+\delta z)\approx \phi(z_0)+J_\phi(z_0)\delta z.
\]
If the observed activity is
\[
x=\phi(z_0)+\delta x,
\]
and \(\delta x\) is small and mostly tangent to the manifold, then a least-squares estimate of the coordinate perturbation solves
\[
\min_{\delta z}\|\delta x-J_\phi(z_0)\delta z\|_2^2.
\]
When \(J_\phi^\top J_\phi\) is invertible,
\[
\boxed{
\widehat{\delta z}
=(J_\phi^\top J_\phi)^{-1}J_\phi^\top\delta x.
}
\]
This is local decoding from neural activity to manifold coordinates. It reuses the projection and least-squares machinery from Chapters 5 and 7.

\paragraph{Worked examples.}
\emph{Small symbolic example.}
A ring manifold in two observed coordinates is
\[
\phi(\theta)=(\cos\theta,\sin\theta),
\qquad \theta\in[0,2\pi).
\]
Its tangent vector is
\[
J_\phi(\theta)=
\begin{bmatrix}
-\sin\theta\\
\cos\theta
\end{bmatrix}.
\]
At \(\theta=0\), the point is \((1,0)\) and the tangent direction is \((0,1)\). Small changes in angle move activity upward along the circle.

\emph{Computational neuroscience example.}
Head-direction cells can form a ring-like population code where each position on the ring corresponds to a direction. Motor-cortex activity during repeated reaching can form condition-dependent trajectories whose low-dimensional coordinates predict movement variables. In both cases, the manifold is an inferred structure in population activity, not a directly visible anatomical object.

\emph{AI example.}
Generative models often learn latent manifolds where nearby latent coordinates decode to similar images, sounds, or actions. Representation-learning systems may align semantic variables with low-dimensional directions or surfaces. The same caution applies as in neuroscience: a visualization of embeddings suggests structure, but quantitative tests are needed to support manifold claims.

\paragraph{Common misconceptions and failure modes.}
Not every low-dimensional projection is a manifold. A manifold claim requires local coordinates and local tangent structure, at least approximately. Intrinsic dimension is not the same as the number of plotted axes. A manifold can curve, fold, or self-approach in a projection. Behavioral correlation along a manifold does not prove that the manifold causes behavior; perturbations and causal analysis are needed for stronger claims.

\paragraph{Exercises.}
\begin{enumerate}[label=\textbf{48.4.\arabic*.}, leftmargin=*]
\item \textbf{Mechanical.} For \(\phi(z)=(z,z^2,z^3)\), with \(z\in\mathbb R\), compute \(J_\phi(z)\) and state the intrinsic dimension of the curve.
\item \textbf{Proof or derivation.} Derive the local least-squares coordinate estimate
\[
\widehat{\delta z}=(J^\top J)^{-1}J^\top\delta x.
\]
\item \textbf{Numerical/computational.} For the ring map \(\phi(\theta)=(\cos\theta,\sin\theta)\), compute \(\phi(0)\), \(\phi(\pi/2)\), and \(\phi(\pi)\). Sketch the corresponding path.
\item \textbf{Interpretation.} What evidence would distinguish a genuine low-dimensional neural manifold from a 2D projection artifact?
\item \textbf{Debug the misconception.} A student says, "If behavior is decodable from a manifold coordinate, that coordinate must be the brain's chosen variable." Explain why decoding is weaker than a causal or mechanistic claim.
\end{enumerate}

\paragraph{Closing bridge.}
Population geometry turns neural recordings into points, clouds, trajectories, latent states, and manifolds. Chapter 49 asks how recurrent circuits can generate such structures. Fixed points become attractors, recurrent weights shape stability, oscillations organize trajectories, and mean-field or neural-field models reduce many interacting neurons to tractable population dynamics.

\clearpage
